\section{Pruebas de validación}

Las pruebas de validación tienen por objetivo verificar el correcto funcionamiento del
sistema en los dos escenarios de despliegue y caracterizar su comportamiento en términos
de latencia y fiabilidad del registro de telemetría.

\subsection{Configuración del entorno de pruebas}

Para las pruebas se emplea la siguiente configuración:
\begin{itemize}
  \item Cuatro fuentes de vídeo simuladas mediante scripts FFmpeg (barras de color con
        reloj en tiempo real a 1920×1080 @25~fps).
  \item Fuentes auxiliares: break, intro, stream blanker y audio de fondo.
  \item Interfaz gráfica voctogui conectada desde un segundo equipo (Escenario 1) o
        desde el host anfitrión (Escenario 2).
\end{itemize}

% --- Tabla sugerida ---
% \begin{table}[H]
%   \centering
%   \caption{Especificaciones hardware de los equipos empleados en las pruebas.}
%   \label{tab:hardware_pruebas}
%   \begin{tabular}{|l|l|l|}
%     \hline
%     \textbf{Equipo} & \textbf{CPU} & \textbf{RAM / Red} \\
%     \hline
%     PC Servidor (voctocore)  & Intel Core i7-xxxx @ x núcleos & xx GB / 1 GbE \\
%     PC Realizador (voctogui) & Intel Core i5-xxxx @ x núcleos & xx GB / 1 GbE \\
%     \hline
%   \end{tabular}
% \end{table}

\subsection{Medición de latencia de conmutación}

Una de las métricas más relevantes para un sistema de producción en directo es la
latencia de conmutación: el tiempo que transcurre desde que el realizador pulsa un
botón de cambio de fuente o composite hasta que el cambio es efectivo y visible en el
monitor de programa.

% TODO: incluir medición real de la latencia (cronómetro en fotograma, marca de tiempo
% del comando vs. evento de telemetría recibido).
% Resultado esperado: latencia < 1 fotograma (< 40 ms) en escenario local.

\begin{itemize}
  \item \textbf{Método de medición:} se registra la marca de tiempo del comando de
        control enviado por la GUI (evento en el log del protocolo TCP) y la marca de
        tiempo del primer fotograma con el nuevo estado visible en el monitor. La
        diferencia entre ambas marcas constituye la latencia de conmutación.
  \item \textbf{Resultado esperado:} latencia inferior a un fotograma a 25~fps
        ($<$40~ms) en el escenario de dos PCs en red local.
\end{itemize}

\subsection{Verificación del sistema de telemetría en RabbitMQ}

Para verificar que todos los eventos de la sesión quedan correctamente registrados en
RabbitMQ, se realiza el siguiente procedimiento:

\begin{enumerate}
  \item Se lanza el sistema completo en el Escenario 2 (Docker Compose).
  \item Se ejecuta una secuencia de acciones representativas: cambio de fuente,
        cambio de composite, activación de stream blanker PAUSE y LIVE, carga y
        descarga de overlay.
  \item Se accede a la consola de administración de RabbitMQ
        (\texttt{http://localhost:15672}) y se comprueba que la cola contiene un
        mensaje por cada acción realizada.
  \item Se ejecuta el script \texttt{rabbit\_consumer\_example.py} para visualizar
        el contenido de los mensajes en formato JSON.
\end{enumerate}

% --- Figura sugerida ---
% \begin{figure}[H]
%   \centering
%   \includegraphics[width=0.85\textwidth]{figuras/rabbitmq_cola.png}
%   \caption{Consola de administración de RabbitMQ mostrando los mensajes de telemetría
%            registrados durante una sesión de prueba.}
%   \label{fig:rabbitmq_cola}
% \end{figure}

% --- Figura sugerida ---
% \begin{figure}[H]
%   \centering
%   \includegraphics[width=0.85\textwidth]{figuras/gui_en_funcionamiento.png}
%   \caption{Captura de pantalla de voctogui durante una sesión de prueba con cuatro
%            fuentes activas, composite pip y overlay cargado.}
%   \label{fig:gui_prueba}
% \end{figure}

\subsection{Resultados}

% TODO: completar con los resultados reales de las pruebas.

Los resultados de las pruebas confirman el correcto funcionamiento del sistema en
ambos escenarios de despliegue. La latencia de conmutación medida es inferior a un
fotograma en el escenario de red local, lo que proporciona una experiencia de
realización fluida sin retardos perceptibles. El sistema de telemetría registra
correctamente todos los eventos de sesión en RabbitMQ, con la marca de tiempo y
el contenido esperados en cada mensaje.
